\medskip\noindent\textbf{Part 1 (15 points): Einstein Ring}

In February 2025, the Euclid mission reported an Einstein ring around the
nearby elliptical galaxy NGC 6505, located at a distance $d_L$ from Earth
(Figure T09.1). The ring is formed when light from a source galaxy at a
distance $d_S \gg d_L$ passes near NGC 6505 and is bent by its gravitational
field.

\ioaafig{2026_TH_T09-f1.pdf}

\begin{description}
\item[(a)] (2 points) The measured redshifts of the source and lens galaxies
  are $0.406$ and $0.042$, respectively. Calculate $d_S$ and $d_L$ (in
  $\mathrm{Mpc}$) assuming that the Newtonian approximation is valid and
  distance is proportional to redshift.
\item[(b)] (3 points) In the linear approximation, as shown in Figure T09.2,
  the total deflection angle of a light ray emitted by the source and reaching
  the observer is given by $2\alpha = 2R^*/b$, where $R^*$ is the
  Schwarzschild radius of the lens and $b$ is the distance of closest approach
  of the light ray to the center of the lens. Express the angular radius
  $\theta$ of the Einstein ring as a function of $R^*$, $d_S$ and $d_L$.

  \ioaafig{2026_TH_T09-f2.pdf}
  \ioaafig{2026_TH_T09-f3.pdf}

\item[(c)] (3 points) Euclid measures the angular radius of the ring around
  NGC 6505 to be $2.5\,\mathrm{arcsec}$. Determine the distance of closest
  approach $b$ (in $\mathrm{parsecs}$) and the mass of the lens galaxy $M$ (in
  units of $\mathrm{M}_\odot$).
\item[(d)] (7 points) Suppose that the lens is displaced by a small distance
  $\varepsilon$ from the line joining the source and the observer
  ($\varepsilon \ll b$) as illustrated in Figure T09.3. Determine the
  quantities $x$ and $x'$ (indicated in the figure), as well as the angular
  displacement $\Delta\theta$ of the Einstein ring image.
\end{description}

\medskip\noindent\textbf{Part 2 (5 points): Gravitational Delay}

Let's now consider light being lensed by the Sun.

\begin{description}
\item[(a)] (2 points) Calculate the total deflection angle $2\alpha$ (in
  $\mathrm{arcsec}$) of a light ray grazing the Sun's surface.
\item[(b)] (3 points) A light ray propagating in the gravitational field of a
  massive object experiences an additional travel time, known as the
  gravitational delay ($\delta$).

  \ioaafig{2026_TH_T09-f4.pdf}

  Consider two light rays passing near the Sun with distances of closest
  approach ($R$) differing by a small amount $\Delta R$ (Figure T09.4). The ray
  with the larger distance of closest approach travels an additional path
  length of approximately $2\alpha\Delta R$, resulting in a difference in
  gravitational delay $\Delta\delta \sim 2\alpha\Delta R/c$. Therefore, the
  absolute value of the derivative of the gravitational delay is
  \[ \frac{d\delta}{dR} = \frac{2\alpha}{c} \]
  Determine the gravitational delay for a light signal from a distant source
  passing near the Sun's surface and reaching the Earth.
\end{description}
